% GENERATED FILE. Edit canonical lessons, then run npm run generate:books.
% canonical-source-hash: fnv1a64:8eccf497c9fd2821
% canonical-lessons: ES-C533-casualidad, ES-C533-coincidencia, ES-C533-fortuito, ES-C533-inesperado, ES-C533-repentino

\chapter{Por casualidad}
\label{ch:es-por-casualidad}
\hlchaptermodality{533}

\begin{chapteropening}
\textbf{By the end of this chapter:} \emph{I can say that something happened by chance, by coincidence or by accident, when nobody expected it.}

Complete the last lesson of chapter 533: I can say that something happened by chance, by coincidence or by accident, when nobody expected it.
\end{chapteropening}

\section[la casualidad]{la casualidad — chance, coincidence}

\label{lesson:ES-C533-casualidad}



\begin{quote}
Before the new one: say the Spanish for to hide, then the Spanish for the secret.
\end{quote}

\subsection*{You'll want to know: la casualidad}
\textbf{la casualidad} — \textquotedblleft{}chance, coincidence\textquotedblright{}.

\begin{quote}
Nos vimos por casualidad en el mercado.
\end{quote}

\begin{quote}
\emph{We met by chance at the market.}
\end{quote}

\emph{Por casualidad} — by chance. \emph{¡Qué casualidad!} — what a coincidence!

\begin{cousinweb}
\emph{casualidad} is built on \emph{casual} 'happening by chance', from Latin \emph{cāsuālis}, from \emph{cāsus} 'a fall, a chance' (Lewis \& Short), and \emph{cāsus} is built on \emph{cadere} 'to fall'.
\end{cousinweb}

\begin{grammarlens}[title={What you've built}]
One more word for telling, explaining and describing at B1.
\end{grammarlens}

\subsection*{Your turn}
\begin{itemize}
\raggedright
  \item \emph{Say it:} \emph{la casualidad}
  \item \emph{Say it:} \emph{la casualidad}, once more
  \item \emph{Say it:} a sentence of your own with \emph{la casualidad}
  \item \emph{Recall:} say the Spanish for to hide, then the Spanish for the secret, then say \emph{la casualidad} again
\end{itemize}

\begin{tcolorbox}[breakable,colback=teal!4,colframe=teal!35!black,title={Before you move on}]
What does la casualidad mean? (\textquotedblleft{}Chance, coincidence\textquotedblright{}.) Say it once more.
\end{tcolorbox}
\section[la coincidencia]{la coincidencia — the coincidence}

\label{lesson:ES-C533-coincidencia}



\begin{quote}
Before the new one: say the Spanish for the secret, then the Spanish for chance.
\end{quote}

\subsection*{You'll want to know: la coincidencia}
\textbf{la coincidencia} — \textquotedblleft{}the coincidence\textquotedblright{}.

\begin{quote}
No fue una coincidencia.
\end{quote}

\begin{quote}
\emph{It was not a coincidence.}
\end{quote}

\emph{Coincidir} is to coincide or to agree.

\begin{cousinweb}
\emph{coincidencia} is built on \emph{coincidir}, from Medieval Latin \emph{coincidere} 'to fall together', usually explained as \emph{co-} and \emph{incidere} 'to fall upon' (Lewis \& Short), built on \emph{cadere} 'to fall'.
\end{cousinweb}

\begin{grammarlens}[title={What you've built}]
One more word for telling, explaining and describing at B1.
\end{grammarlens}

\subsection*{Your turn}
\begin{itemize}
\raggedright
  \item \emph{Say it:} \emph{la coincidencia}
  \item \emph{Say it:} \emph{la coincidencia}, once more
  \item \emph{Say it:} a sentence of your own with \emph{la coincidencia}
  \item \emph{Recall:} say the Spanish for the secret, then the Spanish for chance, then say \emph{la coincidencia} again
\end{itemize}

\begin{tcolorbox}[breakable,colback=teal!4,colframe=teal!35!black,title={Before you move on}]
What does la coincidencia mean? (\textquotedblleft{}The coincidence\textquotedblright{}.) Say it once more.
\end{tcolorbox}
\section[fortuito]{fortuito — accidental, by chance}

\label{lesson:ES-C533-fortuito}



\begin{quote}
Before the new one: say the Spanish for chance, then the Spanish for the coincidence.
\end{quote}

\subsection*{You'll want to know: fortuito}
\textbf{fortuito} — \textquotedblleft{}accidental, by chance\textquotedblright{}.

\begin{quote}
Fue un hallazgo fortuito: buscaba otra cosa.
\end{quote}

\begin{quote}
\emph{It was a chance find: I was looking for something else.}
\end{quote}

\emph{Un caso fortuito} is an accident nobody could foresee.

\begin{cousinweb}
\emph{fortuito} is a learned borrowing from Latin \emph{fortuītus} 'happening by chance' (EtymDB), built on \emph{fors} 'chance, luck' (Lewis \& Short).
\end{cousinweb}

\begin{grammarlens}[title={What you've built}]
One more word for telling, explaining and describing at B1.
\end{grammarlens}

\subsection*{Your turn}
\begin{itemize}
\raggedright
  \item \emph{Say it:} \emph{fortuito}
  \item \emph{Say it:} \emph{fortuito}, once more
  \item \emph{Say it:} a sentence of your own with \emph{fortuito}
  \item \emph{Recall:} say the Spanish for chance, then the Spanish for the coincidence, then say \emph{fortuito} again
\end{itemize}

\begin{tcolorbox}[breakable,colback=teal!4,colframe=teal!35!black,title={Before you move on}]
What does fortuito mean? (\textquotedblleft{}Accidental, by chance\textquotedblright{}.) Say it once more.
\end{tcolorbox}
\section[inesperado]{inesperado — unexpected}

\label{lesson:ES-C533-inesperado}



\begin{quote}
Before the new one: say the Spanish for the coincidence, then the Spanish for accidental.
\end{quote}

\subsection*{You'll want to know: inesperado}
\textbf{inesperado} — \textquotedblleft{}unexpected\textquotedblright{}.

\begin{quote}
Recibimos una visita inesperada.
\end{quote}

\begin{quote}
\emph{We had an unexpected visit.}
\end{quote}

\begin{cousinweb}
\emph{inesperado} is built on \emph{esperado}, the participle of \emph{esperar}, with \emph{in-} 'not'. \emph{Esperar} is inherited from Latin \emph{spērāre} 'to hope' (Lewis \& Short).
\end{cousinweb}

\begin{grammarlens}[title={What you've built}]
One more word for telling, explaining and describing at B1.
\end{grammarlens}

\subsection*{Your turn}
\begin{itemize}
\raggedright
  \item \emph{Say it:} \emph{inesperado}
  \item \emph{Say it:} \emph{inesperado}, once more
  \item \emph{Say it:} a sentence of your own with \emph{inesperado}
  \item \emph{Recall:} say the Spanish for the coincidence, then the Spanish for accidental, then say \emph{inesperado} again
\end{itemize}

\begin{tcolorbox}[breakable,colback=teal!4,colframe=teal!35!black,title={Before you move on}]
What does inesperado mean? (\textquotedblleft{}Unexpected\textquotedblright{}.) Say it once more.
\end{tcolorbox}
\section[repentino]{repentino — sudden}

\label{lesson:ES-C533-repentino}



\begin{quote}
Before the new one: say the Spanish for accidental, then the Spanish for unexpected.
\end{quote}

\subsection*{You'll want to know: repentino}
\textbf{repentino} — \textquotedblleft{}sudden\textquotedblright{}.

\begin{quote}
El cambio de tiempo fue repentino.
\end{quote}

\begin{quote}
\emph{The change in the weather was sudden.}
\end{quote}

\emph{De repente} — suddenly — comes from the same family.

\begin{cousinweb}
\emph{repentino} is a learned borrowing from Latin \emph{repentīnus} 'sudden' (EtymDB), built on \emph{repēns} 'sudden, unexpected'.
\end{cousinweb}

\begin{grammarlens}[title={What you've built}]
One more word for telling, explaining and describing at B1.
\end{grammarlens}

\subsection*{Your turn}
\begin{itemize}
\raggedright
  \item \emph{Say it:} \emph{repentino}
  \item \emph{Say it:} \emph{repentino}, once more
  \item \emph{Say it:} a sentence of your own with \emph{repentino}
  \item \emph{Recall:} say the Spanish for accidental, then the Spanish for unexpected, then say \emph{repentino} again
\end{itemize}

\begin{tcolorbox}[breakable,colback=teal!4,colframe=teal!35!black,title={Before you move on}]
What does repentino mean? (\textquotedblleft{}Sudden\textquotedblright{}.) Say it once more.
\end{tcolorbox}
